\section{Cross-validation and checkpoint transport}
\label{sec:validation}

Two independent checks constrain the method externally: a cross-validation
against PIConGPU's exact-quadrature field reference, and a checkpoint and
restart transport. Neither is a truth comparison.

\subsection{Cross-validation against the PIConGPU reference}

The reference is the TSC-convolved field (Section~\ref{sec:accuracy}),
evaluated in PIConGPU~\citep{burau2010} in its reference-only
field-initializer mode. The comparison passes for every sampled point, and
both sides use identical Yee offsets, grid, and origin; the reference
staggered point maps to a picBunch node by index with zero snap deviation.

At $318$ staggered points on a $64\times64\times32$ grid, cell
$10^{-7}$\,m, the overall worst per-component relative error is
$4.86\times10^{-2}$ for $\gamma=1$ ($E/z$) and $6.20\times10^{-1}$ for
$\gamma=10$ ($E/z$; $E/y$ is $5.16\times10^{-1}$), against a reference
reproducibility floor of order $10^{-6}$ (reported from the
\texttt{rosi\_v100} campaign, original provenance the PIConGPU
\texttt{referenceOnly} run; the comparison file is archived outside this
repository and was not re-measured in this campaign). picBunch is about four
to six orders of
magnitude above that floor; the measured error is concentrated at the core
and the beamed peaks, where the TSC support is, while the bulk agrees to
below $1$--$3\%$ and $B$ is exactly zero for $\gamma=1$ in both components.

The whole-field comparison, $N=1000$ on $32\times32\times16$, cell
$2\times10^{-7}$\,m, $16384$ nodes per component, gives the same picture:
$\gamma=1$ bulk median $3.8$--$7.8\times10^{-3}$ with an L2 of
$4.6$--$5.7\times10^{-2}$; $\gamma=10$ bulk median $1.4$--$3.3\times10^{-2}$
but a beamed-core overshoot that pushes the $E/y$, $E/z$ L2 to
$0.76$--$0.92$. The forward $E/x$ is well resolved at $\gamma=10$
($\sim6\%$); the beamed transverse $E/y$, $E/z$ are $50$--$130\%$ worst at
the core, where the Lorentz-contracted field is a sub-cell pancake below
the TSC support and the spectral near-field kernel resolution.

\begin{table}[t]
\centering
\small
\caption{Whole-field comparison of picBunch against the PIConGPU
\texttt{referenceOnly} field, $N=1000$, $32\times32\times16$, cell
$2\times10^{-7}$\,m, $16384$ nodes/component, rosi\_v100.
\texttt{worst\_rel} is peak-normalized; median and p99 are per-point over
nodes with $|\mathrm{ref}|>10^{-3}\times$peak; L2 is pooled. Not a truth
comparison.}
\label{tab:wholefield}
\begin{tabular}{lrrrr}
\toprule
component & worst\_rel & median & p99 & L2 \\
\midrule
\multicolumn{5}{l}{\emph{$\gamma=1$} (B identically 0)} \\
$E/x$ & $1.48\times10^{-1}$ & $3.92\times10^{-3}$ & $3.20\times10^{-1}$ & $4.62\times10^{-2}$ \\
$E/y$ & $1.49\times10^{-1}$ & $3.76\times10^{-3}$ & $2.93\times10^{-1}$ & $4.75\times10^{-2}$ \\
$E/z$ & $1.22\times10^{-1}$ & $7.81\times10^{-3}$ & $3.65\times10^{-1}$ & $5.67\times10^{-2}$ \\
\midrule
\multicolumn{5}{l}{\emph{$\gamma=10$} (B/x identically 0)} \\
$E/x$ & $1.36\times10^{-1}$ & $4.80\times10^{-3}$ & $1.23\times10^{-1}$ & $4.44\times10^{-2}$ \\
$E/y$ & $1.98$  & $1.36\times10^{-2}$ & $7.25$ & $7.63\times10^{-1}$ \\
$E/z$ & $1.93$  & $3.33\times10^{-2}$ & $10.1$ & $9.19\times10^{-1}$ \\
$B/y$ & $4.42\times10^{-1}$ & $3.12\times10^{-2}$ & $1.41$ & $2.58\times10^{-1}$ \\
$B/z$ & $4.63\times10^{-1}$ & $1.52\times10^{-2}$ & $0.93$ & $2.20\times10^{-1}$ \\
\bottomrule
\end{tabular}
\end{table}

The disagreement is per-source, local to the TSC support, and independent
of particle count: on a $64\times64\times32$, cell $10^{-7}$\,m, stationary
bunch ($\gamma=1$) the band is flat in $N$ from $1$ to $1000$
($10^{-3}$--$3\times10^{-2}$). The $O(1)$ values at $N\ge10$ occur only
because the global worst sample point falls $0.415$--$0.65$ cells from a
source, inside that source's TSC support; an incorrect charge weighting would
grow with $N$, and it does not.

\subsection{Checkpoint format and bitwise arrival}

The checkpoint writer emits a PIConGPU openPMD field checkpoint matching
PIConGPU's own plugin schema: root layout, openPMD standard version 1.1.0
\citep{openpmd} (written by the openPMD-api 0.17.1 software), $E$, $B$
$x/y/z$ at the exact Yee positions, grid spacing
and offset, and the \texttt{picongpu\_idProvider} mesh metadata. The
default precision is uniform \texttt{float32} for every field and particle
record, so a stock single-precision PIConGPU can load the file;
\texttt{float64} is a uniform opt-in that records a warning attribute
because stock PIConGPU cannot load double records. Every species from the
input bunch is copied, the \texttt{charge} and \texttt{mass} records are
written rather than omitted, and the species is a single particle patch,
\texttt{patch} = \texttt{"bbox"} or \texttt{"domain"}. For a multi-rank
restart the writer splits the species into one patch per output rank:
\texttt{patch\_count} in \texttt{[placement]} gives cumulative, disjoint
\texttt{numParticlesOffset} ranges, with one-cell overlap between
neighbouring $z$ regions, and every particle belongs to exactly one patch.

Field arrival is bitwise. A fresh checkpoint written from a two-species
input (16 electrons and 16 ions, \texttt{picongpu\_idProvider} startId $0$,
nextId $32$) restarts in PIConGPU with exit 0 under both \texttt{patch}
modes (\texttt{bbox}: offset $(7,8,3)$, extent $(8,10,6)$;
\texttt{domain}: offset $(0,0,0)$, extent $(24,24,12)$), all six field
components bitwise-equal, $\max_\mathrm{abs\,diff} = 0.0$, electrons with
ids $0$ to $15$, first-step displacement exactly $\beta c\,dt$ with
$\beta = 0.9949874367$. The multi-rank restart is verified on the
$24\times24\times96$ grid with the $4$ ranks in $z$ ($-d\,1\,1\,4$,
per-rank $24\times24\times24$), which is an integer multiple of the
per-rank grid on every axis and so needs no automatic grid adjustment:
a $256$-particle bunch written with \texttt{patch\_count} $4$ (four
particle patches, one per rank, with two empty patches for this compact
bunch) restarts PIConGPU with exit 0, and after one step all six field
components are bitwise-equal to the checkpoint,
$\max_\mathrm{abs\,diff}=0.0$
\citep[measured in]{restart4measure}.

\paragraph{Placement and particle-data integrity.}
A 12-marker bunch with an asymmetric position and momentum pattern was
placed at the xz middle and the y first quarter of a $32\times64\times32$
domain.  picBunch reports centroid fractions
$(0.500000,0.250000,0.500000)$; after one PIConGPU step the centroid is
$(0.500755,0.250446,0.501466)$, the raw momentum records are preserved
exactly, and the per-particle displacement matches $dt\,c\,u/\gamma$ to
$4.1\times10^{-7}$ cells.  The displacement-to-momentum ratio is one
positive scalar on all three axes (spread $8.9\times10^{-6}$), so the
restart has no axis permutation and no mirror.  One metadata requirement is
sharp: PIConGPU restart ignores the particle \texttt{unitSI} and interprets
momentum in \texttt{sim.unit.mass()}$\cdot$\texttt{unit\_speed}, so a
checkpoint written with the electron mass as \texttt{unit\_mass} restarts
with all momenta $10^{4}$ too large; \texttt{unit\_mass} must equal the
deck's \texttt{sim.unit.mass()}.

On a $24\times24\times12$ deck the restart is exit 0 with
\texttt{--periodic 1 1 1}, and the $16$ electrons load; initialization
drops from the $\sim116$\,s neuralBunch path to $1.8$\,s because the
checkpoint carries the field. For an open-boundary deck the writer copies
the \texttt{Convolutional PML} slabs from the input checkpoint verbatim
(converting a non-float32 slab to float32), so the PML state arrives with
the restart. On a real LWFA checkpoint ($638976$ electrons, twelve
\texttt{Convolutional PML} groups with six datasets each,
$64\times128\times64$, \texttt{--periodic 0 0 1}) the writer copies all
twelve groups and the restart exits 0; the copied PML datasets are float32
and bitwise equal to the source, initialization is $3.986$\,s, and the
main $E$ and $B$ records are recomputed from the particles and stay
float32.
